\originalchapter{高斯定律与对称性}\label{ch:gauss}
\originalsection{通量与包围电荷}
点电荷在以它为球心的半径$r$球面上给出通量$E4\pi r^2=q/\eps$. 对任意闭合面, 一个小面元的径向投影面积除以$r^2$就是从电荷看去的立体角. 当电荷在面内时所有小元立体角合为$4\pi$, 在面外时穿入和穿出抵消. 结合叠加即得高斯定律:
\begin{equation}\label{eq:gauss}
\oint_S\ee\cdot\dd\vct A=\frac{Q_{\rm enc}}\eps.
\end{equation}
左边是整个闭合面的净通量, 右边只含面内电荷. 面外电荷通常仍影响面上每一点的场, 只是对净通量贡献为零. 因此$Q_{\rm enc}=0$不推出$\ee=0$.

高斯定律对任意形状成立, 但求未知场时还需对称性, 使场在所选曲面上等大且方向已知. 若只有一个总通量关系而未知场随位置变化, 则不能直接由通量求出局部场.

\originalsection{球、柱和平面对称}
均匀体电荷球半径$R$, 总电荷$Q$. 旋转对称性使场径向且只依赖$r$. 半径$r<R$高斯球内有$Q(r/R)^3$, 于是
\[
\ee(r)=\begin{cases}\displaystyle\frac{Qr}{4\pi\eps R^3}\uv r,&r<R,\\[3pt]
\displaystyle\frac Q{4\pi\eps r^2}\uv r,&r>R.\end{cases}
\]
中心场为零, 边界两侧连续. 对仅在球面上均匀带电的薄壳, 内部无包围电荷并有球对称, 所以内场为零, 外场仍等同于中心点电荷.

无限长直线电荷密度$\lambda$. 轴向平移和绕轴转动对称, 场为$E(s)\uv s$. 取长$\ell$的同轴高斯柱, 两端通量零、侧面为$E2\pi s\ell$, 得$E=\lambda/(2\pi\eps s)$. 有限长导线靠近末端失去该对称性; “足够长”近似须限制观测区远离末端.

无限均匀薄面密度$\sigma$, 两侧场法向且大小相等. 取穿过面的扁柱, 总通量为$2EA$, 得一侧$E=\sigma/(2\eps)$, 正电荷指向离开平面的方向. 一对相反面密度$\pm\sigma$的平行大板, 两板中间场相加为$\sigma/\eps$, 外侧相消. 有限板在边缘有散边场, 匀场模型要求间距远小于横向尺度且避开边缘.

\originalsection{边界与势差}
穿过面电荷的小扁柱给出法向跃变$E_{n,+}-E_{n,-}=\sigma/\eps$. 静电场环量为零, 取跨界小矩形则切向分量连续. 这些边界关系说明面电荷上场可以跳跃, 而普通有限体密度界面通常无这样的法向跳跃.

厚度$2a$、沿$y,z$无限延展的均匀带电板占$|x|<a$, 体密度$\rho$. 取对称扁柱可得内场$E_x=\rho x/\eps$, 外场$E_x=\rho a\operatorname{sgn}(x)/\eps$. 内部势$V(x)-V(0)=-\rho x^2/(2\eps)$. 高斯法、势的导数和边界连续是相互核对的三个入口.
